\subsubsection{Enclosures and faces}
We keep the notations introduced at the end of Section~2.1. Given a filter $F$ on $\jepPubBbb{A}$, we define its enclosure $\operatorname{cl}_{\jepPubBbb{A}}(F)$ as the filter of all subsets of $\jepPubBbb{A}$ containing some element of $F$ of the form $\bigcap_{\alpha\in\Delta}D(\alpha,k_\alpha)$, with $k_\alpha\in\jepPubBbb{Z}\cup\{+\infty\}$ for any $\alpha\in\Delta$. Sets of the form $D(\alpha,k)$ with $\alpha\in\Phi$ and $k\in\jepPubBbb{Z}$ are called half-apartments in $\jepPubBbb{A}$. Sets of the form $M(\alpha,k):=\{t\in\jepPubBbb{A}\mid\alpha(t)+k=0\}$ with $\alpha\in\Phi$ and $k\in\jepPubBbb{Z}$ are called walls in $\jepPubBbb{A}$.

A local face in $\jepPubBbb{A}$ is a filter $F^\ell$ of the form $\operatorname{germ}_x(x+F^v)$, where $x\in\jepPubBbb{A}$ is called the vertex of $F^\ell$ and a vectorial face $F^v\subset\jepPubBbb{A}$ is a vectorial face called the direction of $F^\ell$. To keep track of the elements $x$ and $F^v$, such a local face may be denoted $F^\ell(x,F^v)$. A local face is said to be spherical when its direction is spherical: in this case, its pointwise stabilizer under the action of $W^v$ is a finite group.

A face in $\jepPubBbb{A}$ is a filter $F=F(x,F^v)$ associated with a point $x\in A$ and a vectorial face $F^v\subset A$ as follows: a subset $S$ of $\jepPubBbb{A}$ belongs to $F$ iff it contains an intersection of half-spaces $D(\alpha,k_\alpha)$ or open half-spaces $\mathring{D}(\alpha,k_\alpha)$, with $\alpha\in\Delta$ and $k_\alpha\in\jepPubBbb{Z}$, that also contains the local face $F^\ell(x,F^v)$. Note that (local) faces can be ordered as follows: given two such faces $F,F'$ in $\jepPubBbb{A}$, we say that $F$ is a face of $F'$ (or $F'$ contains $F$, or $F'$ dominates $F$) when $F\subset F'$.

As explained at the end of Section~2.1, the action of $W$ on $\jepPubBbb{A}$ permutes the sets of the form $D(\alpha,k)$, where $(\alpha,k)$ runs over $\Delta\times\jepPubBbb{Z}$. In particular, this implies that $W$ permutes enclosures (resp. walls, faces) of $\jepPubBbb{A}$.

The dimension of a face $F$ is the smallest dimension of an affine space generated by some element $S$ of $F$. Such an affine space is unique and is called the support of $F$. A local chamber (or local alcove) is a maximal local face, i.e., a local face of the form $F^\ell(x,\pm w\cdot C_f^v)$ for $x\in\jepPubBbb{A}$ and $w\in W^v$. The fundamental local chamber is $C_0^+:=F^\ell(0,C_f^v)$. A local panel is a spherical local face which is maximal among faces that are not chambers. Equivalently, a local panel is a spherical local face of dimension $\dim\jepPubBbb{A}-1$. Analogue definitions of chambers and panels exist, see for instance~\cite[\S1.4]{jep88refGR08}. Finally, a local face is of type $0$ (or: is a type $0$ local face) if its vertex lies in $Y$. We denote by $F_0$ the local face $F^\ell(0,\jepPubBbb{A}_{\mathrm{in}})$, where $\jepPubBbb{A}_{\mathrm{in}}=F^v(I)=\bigcap_{i\in I}\ker(\alpha_i)$. From now on, we will write type $0$ face instead of type $0$ local face to make it shorter.
\begin{remark}\label{jep88localRemark2_3}
In~\cite{jep88refRou11}, Rousseau defines a notion of chimney that he uses in his axiomatization of masures. We do not define here what chimneys are: we only recall that each sector-germ is a splayed, solid chimney-germ, that each spherical face is contained in a solid chimney and that the action of $W$ on $\jepPubBbb{A}$ permutes the chimneys and preserve their properties (being splayed or solid for instance). For more details about this, see~\cite[\S1.10]{jep88refRou11}.
\end{remark}
